% =====================================================================
% Hardware Implementation  (~430 words; physical-build sub-sections
% intentionally light per author note)
% =====================================================================
\section{Hardware Implementation}
\label{sec:hardware}

\subsection{Bill of materials and physical assembly}
The arm is mechanically based on the EEZYbotARM family of 3D-printed
designs, with original modifications detailed in
Section~\ref{sec:cad}. The principal electronic components are
listed in Section~\ref{sec:design}; the controller carries two KY-023
analogue joysticks, a momentary mode-switch button and an HC-05
master radio.

% TODO PHOTO: assembled arm on the bench, ideally with cabling tidy
% and the camera module visible mounted to the base. Replace the
% placeholder text inside the figure with \includegraphics{arm_photo}.
\begin{figure}[t]
  \centering
  \fbox{\parbox{0.92\columnwidth}{\centering\vspace{4pt}
    \emph{[Photograph of the assembled arm on the bench, cabling
    visible. To be added once final assembly is complete.]}\vspace{4pt}}}
  \caption{Assembled arm prototype. Power and signal cabling are
  separated; the gripper end-effector is visible at upper right.}
  \label{fig:armphoto}
\end{figure}

% Physical-build write-up to be filled in by group member responsible
% for assembly. Topics to cover (placeholders below):
% - frame assembly steps and torque values
% - servo mounting and offset calibration
% - cable routing and strain relief
% - end-stop verification and physical limits
\subsubsection*{Mechanical assembly (placeholder)}
\emph{This subsection is intentionally short in the first draft and
will be expanded by the group member responsible for the physical
build. Topics to cover include the assembly sequence, torque values
used on the servo horns, cable routing and strain relief, and the
mechanical end-stop verification procedure.}

\subsection{Sensor interfacing}
\label{sec:sensors}
\textbf{Joystick axes.} Each KY-023 is a centred-detent potentiometer
network with 0--3.3\,V output. The Uno R4's 12-bit ADC samples each
axis at $\geq$200\,Hz and applies a small dead-band to suppress
quiescent jitter. The mode-switch button is debounced in firmware
with a 20\,ms window.

\textbf{Joint potentiometers.} Each of the four joints carries a
10\,k$\Omega$ linear potentiometer, geared (or directly coupled) to
the joint axis to provide an absolute position reading. The wiper
voltage is RC-filtered ($R$=1\,k$\Omega$, $C$=100\,nF, 1.6\,kHz cut)
and read by the Mega's 10-bit ADC at 100\,Hz under timer interrupt.
A 16-sample boxcar average removes residual servo-induced noise
without compromising the closed-loop bandwidth.

\textbf{IR proximity (autonomous mode).} A short-range IR module on
a digital input gates the
\texttt{DETECT} state of the autonomous FSM (Section~\ref{sec:fsm}).

\subsection{Wiring schematics}
The full wiring of each unit is shown in
Figs.~\ref{fig:armwiring} and~\ref{fig:ctrlwiring}. The
controller-side HC-05 is voltage-divided on its RX line so that the
Uno's 5\,V TX presents approximately 3.3\,V to the radio
(Fig.~\ref{fig:hc05div}); the same divider is reproduced on the arm
side.

\begin{figure}[t]
  \centering
  \includegraphics[width=0.95\columnwidth]{03_arm_wiring.pdf}
  \caption{Arm-side wiring. The Mega 2560 drives four servos from a
  separate 6\,V rail; the HC-05 slave is on Serial1.}
  \label{fig:armwiring}
\end{figure}

\begin{figure}[t]
  \centering
  \includegraphics[width=0.95\columnwidth]{04_controller_wiring.pdf}
  \caption{Controller-side wiring. Two KY-023 joysticks and a mode
  button feed the Uno R4 Minima; the HC-05 master is on Serial1.}
  \label{fig:ctrlwiring}
\end{figure}

\begin{figure}[t]
  \centering
  \includegraphics[width=0.65\columnwidth]{02_hc05_voltage_divider.pdf}
  \caption{HC-05 RX voltage divider.}
  \label{fig:hc05div}
\end{figure}

\subsection{Power and protection}
The split-rail power architecture is shown in Fig.~\ref{fig:power}.
A common ground reference is established at the Mega; servo $V_{cc}$
is decoupled with a 470\,$\mu$F bulk capacitor at the rail and
100\,nF ceramics at each servo. An e-stop button on the controller
asserts an emergency-stop flag in the next outgoing frame, which the
arm interprets by parking the joints in a safe pose.

\begin{figure}[t]
  \centering
  \includegraphics[width=0.95\columnwidth]{01_power_architecture.pdf}
  \caption{Power architecture. Logic and motor rails are separated;
  grounds are common. The arm's mains rail is never connected to the
  Mega's $V_{in}$.}
  \label{fig:power}
\end{figure}

\corehit{Hardware brief Objectives 1, 3 and~5: hardware components
assembled, sensors interfaced, motor control wired up; circuit
schematics for both transmitter and receiver provided.}

\advhit{Analogue-front-end design (RC filter $+$ digital averaging
on potentiometers); split-rail power architecture with documented
grounding strategy; voltage-level translation on the HC-05 RX line;
hardware e-stop wired through firmware to a fail-safe pose.}
